\chapter{狭义相对论}

\section{相对性原理}

两百多年来，牛顿所阐明的运动方程一直被认为是对自然的一种正确描述。第一次看出这些定律中存在的一个谬误，并且找到了修正它的方法是在1905年，这两件事都是爱因斯坦所提出的。

我们曾用下面的方程表示牛顿第二定律
\begin{equation*}
    F=\frac{d(mv)}{dt}
\end{equation*}
牛顿在叙述这一定律时默认了这样一个假定，即质量$m$是一个恒量，但是我们现在知道这并不正确，而是物体的质量要随着其速度的增加而增大。在经爱因斯坦修正后的公式中，质量具有数值
\begin{equation}
    \label{Eq:I:15:1}
    m=\frac{m_0}{\sqrt{1-v^2/c^2}}
\end{equation}
这里“静止质量”$m$。表示一个不运动的物体的质量，$c$是光速，约为$3 \times 10^8~m \cdot s^{-1}$或$186000~mi \cdot s^{-1}$。

对于那些只想学一点能够解释问题就行了的人来说，这个式子就是全部相对论了——它只是对质量引入一个修正因子来改变一下牛顿定律。从公式本身很容易看出，在通常情况下，质量的增加是十分微小的。即使速度大到像绕地球运转的卫星一样，即约$5~mi \cdot s^{-1}$，于是$v/c-5/186000$；把这个值代入公式后表明，对质量的修正只是20亿到30亿分之一，小得几乎无法观察到。实际上，这个公式的正确性已被对许多种粒子作出的观察所广泛证明，这些粒子运动速度很大，一直到实际上等于光速。然而，由于这种效应通常非常之小，所以看来似乎不寻常的是，它在理论上的发现先于它在实验上的发现。从实验上看，当速度足够大的时候，这种效应非常之大，但是它不是用这一方法发现的。所以，看一下一条涉及到这样一种细致的修正的定律，在它第一次被发现的那个时期如何能从实验和物理推理的结合中得以产生该是十分有趣的。有很多人对这条定律的发现作出了贡献，但是爱因斯坦的发现是这些人的工作的最后成果。

实际上爱因斯坦的相对论包括两部分内容。这一章所谈到的是1905年以来就已存在的狭义相对论。1915年爱因斯坦发表了称为广义相对论的补充理论。这后一个理论所讨论的是将狭义相对论推广到引力定律中去的情况；这里我们将不去讨论它。

相对性原理是牛顿在他的运动定律的一个推论中首先提出的：“封闭在一个给定空间中的诸物体，它们彼此之间的运动是同一的，无论这个空间是处于静止状态还是均匀地沿一直线向前运动。”这意味着，比如说，如果有一艘宇宙飞船在以均匀速度飞行，那么在飞船上所做的所有实验以及所有的现象，将与飞船不运动时所看到的完全相同，当然这是指如果人们并不伸出头去往外看的话。这就是相对性原理的含义。它是一个十分简单的观念，唯一的问题在于是否确实如此：即从一个运动系统内进行的所有实验中得到的物理定律是否看来都与如果该系统处于静止时所得出的相同。让我们首先来研究一下，牛顿定律在运动系统中是否相同。

\begin{wrapfigure}{r}{0.4\textwidth}
    \centering
    \includegraphics[width=0.35\textwidth]{Chapter15/沿着x轴以均匀速度作相对运动的两个坐标系}
    \caption{沿着$x$轴以均匀速度作相对运动的两个坐标系}
    \label{figure:沿着x轴以均匀速度作相对运动的两个坐标系}
\end{wrapfigure}
假设某一个人莫（Moe）以均匀速度沿$x$方向运动，并且测量了某一点的位置，如图\ref{figure:沿着x轴以均匀速度作相对运动的两个坐标系}所示。他把这个点在他这个坐标系中的“$x$距离”记作$x'$。另一个人乔（Joe）静止不动，并且测量了同一点的位置。用他的坐标系中的$x$坐标记作$x$。这两个系统的坐标之间的关系可以从图中清楚看出。经过时间$t$后，莫的原点移动了一段距离$ut$，如果这两个系统原先是重合在一起的，那么
\begin{equation}
    \label{Eq:I:15:2}
    \begin{aligned}
        x' &= x - ut \\
        y' &= y \\
        z' &= z \\
        t' &= t
    \end{aligned}
\end{equation}

如果把这个坐标变换代入牛顿定律中去，那么我们发现，这些定律变换到带撇（$'$）的坐标系中时，仍旧是同样的定律；这就是说，牛顿定律在运动系统中与在静止系统中一样具有相同的形式，因此，依靠力学实验不可能说出系统究竟是否在运动。

相对性原理在力学中已应用了很长一段时间。曾经有许多人，特别是惠更斯（Huygens）,应用它来求出弹子球碰撞的规则，这与我们在第10章中用它来讨论动量守恒时所用的方法很相同。在上一世纪中，由于对电、磁以及光等现象的研究，人们对于这条原理的兴趣更加浓厚了。许多人对这些现象所作的一系列精心研究，其结晶就是麦克斯韦电磁场方程组，它们在统一的体系下描写了电、磁与光的现象。但是麦克斯韦方程组似乎并不遵循相对性原理。这就是说，如果我们用式（\ref{Eq:I:15:2}）代入麦克斯韦方程组并对它进行变换，那么它们的形式不再保持相同；因此，在飞行的宇宙飞船中，电与光的现象应当与飞船静止时不同。这样，我们就可以利用这些光的现象来确定飞船的速度；特别是可以通过适当的光学或电学测量来确定飞船的绝对速度。麦克斯韦方程组的结论之一是，如果在电场中产生扰动，以致有光发射出来，那么这些电磁波在所有方向上均等地而且以相同的速度$c$（即$186000~mi \cdot s$）传播出去。方程组的另一个结论是，假如扰动源在运动，那么所发射的光将以同样的速度$c$穿过空间。这与声的情况相似，声波的速度也与声源的运动无关。

在光的情况下，这种与声源运动的无关性引起了一个有趣的问题。

假定我们以速度驱车前进，从后面射来的光，以速度$c$追过了我们的车子。对式（\ref{Eq:I:15:2}）中的第一个方程微分就得到
\begin{equation*}
    \frac{dx'}{dt}=\frac{dx}{dt} - u
\end{equation*}
这意味着，按照伽利略变换，我们在汽车里测得的掠车而过的光的表观速度不应当是$c$而应

当是（$c-u$）。例如，如果汽车以$100000~mi \cdot s^{-1}$速度前进，而光速是$186000~mi \cdot s^{-1}$，那么
掠车而过的光的表观速度应当是$86000~mi \cdot s^{-1}$。总之，在任何情况下，只要测出掠车而过的光的速度（如果伽利略变换对于光是正确的话），我们就应当可以决定汽车的速度。在这种一般设想的基础上，曾经进行了大量的实验以确定地球的速度，但是它们全都失败了——根本没有发现地球有什么速度。我们将详细地讨论其中的一个实验，用以确切地说明我们做了一些什么以及问题何在；显然，是有一些问题，物理方程有点不对。那么怎么会这样的呢？

\section{洛伦兹变换}

当物理方程在上述情况下的失效暴露出来的时候，所出现的第一个想法就是认为这个麻烦的根源必定在于当时只有20年之久的新的麦克斯韦电动力学方程组。看来相当明显的是，这些方程一定是错误的，所以要做的事就是这样来改变它们，使得相对性原理在伽利略变换下能够得到满足。在这种尝试中，必须在方程组中引入一些新的项，而这些项预言了一些新的电现象，但一旦用实验来检验它们时，这些现象就根本不存在，因而，这个尝试必须予以放弃。于是人们逐渐明白，麦克斯韦电动力学方程组是正确的，必须到其他地方去寻找症结所在。

在这期间，洛伦兹注意到一件令人注目的奇怪的事，那就是当他在麦克斯韦方程组中进行以下代换时
\begin{equation}
    \label{Eq:I:15:3}
    \begin{aligned}
        x'&=\frac{x-ut}{\sqrt{1-\frac{u^2}{c^2}}} \\
        y'&=y \\
        z'&=z \\
        t'&=\frac{t-\frac{ux}{c^2}}{\sqrt{1-\frac{u^2}{c^2}}}
    \end{aligned}
\end{equation}
发现这些方程组在这种变换下保持其原有形式不变！式（\ref{Eq:I:15:3}）现在通称为洛伦兹变换。爱因斯坦仿效原来由庞加莱（Poincaré）提出的设想，作出了这样的一个假设：所有的物理定律都应该是这样的定律，它们在洛伦兹变换下保持不变。换句话说，需要予以改变的，不应当是那些电动力学定律，而应当是力学定律。那么，我们将如何改变牛顿定律，使它们在洛伦兹变换下保持不变呢？如果确定的是这样一个目标，那么我们必须把牛顿方程这样来予以改写，使得它们能满足我们所提出的条件。结果发现，这里的唯一要求，就是牛顿方程中的质量$m$必须代之以等式（\ref{Eq:I:15:1}）中所示的形式。作了这个改变之后，牛顿定律与电动力学定律就会完全协调。这时，如果我们用洛伦兹变换把莫的测量与乔的测量相比较，那么就根本不可能发现究竟谁在运动，因为两个坐标系中，所有方程的形式都是相同的！

当我们用新的时间与坐标之间的变换来代替旧的变换时，这究竟意味着什么，讨论一下这个问题是颇为有趣的，因为旧的变换（伽利略变换）似乎是不证自明的，而新的变换（洛伦兹变换）看来是奇特的。我们希望知道的是，在逻辑上以及在实验上是否可能把新的而不是旧的变换看作是正确的。要弄清楚这一点，单去研究力学定律是不够的，而应像爱因斯坦所做的那样，必须对我们关于时间和空间的观念进行分析，以求得对这种变换的理解。我们将不得不稍微花一点时间来讨论这些观念以及它们对力学的含义，所以我们要说明在前，由于其结果与实验相符合，这种努力是完全有理由的。

\section{迈克耳逊-莫雷实验}

如上所述，人们曾经做过多次尝试，以确定地球通过一种假设的“以太”时的绝对速度，而以太是被想象为充满整个空间的。这些实验中最著名的一个是1887年迈克耳逊
\begin{wrapfigure}{r}{0.5\textwidth}
    \centering
    \includegraphics[width=0.48\textwidth]{Chapter15/迈克尔逊-莫雷实验的示意图}
    \caption{迈克尔逊-莫雷实验的示意图}
    \label{figure:迈克尔逊-莫雷实验的示意图}
\end{wrapfigure}
（Michelson）和莫雷（Morley）所做的。这个实验所得到的负结果经过18年之后才最终由爱因斯坦作出了解释。

迈克耳逊-莫雷实验使用了如图\ref{figure:迈克尔逊-莫雷实验的示意图}所示的一种装置。它主要包括一个光源A，一块部分镀银的玻璃片B，两面镜子C和E。所有这些都装在一个牢固的底座上。两面镜子放在离B都等于L的地方。B片将射来的光分为两束，这两束光以互相垂直的方向分别向两面镜子射去，并在那里被反射而回到B。在返回到B后，这两束光线又作为叠加分量D与F组合起来。如果光线从B到E一次来回的时间与光从B到C一次来回的时间相同，那么所产生的两条光线D与F的相位相同，因而彼此加强。但是如果这两个时间稍有差异，那么两条光线之间就会有一点相位差，结果将产生干涉现象。如果这个装置在以太中“静止”不动，那么这两个时间应该精确相等，但是如果它以速度$u$向右运动，那么这两个时间就应有所差别。让我们看看其原因何在。

首先，我们来计算光从B到E而后再返回到B所需的时间。假设光从B片到镜E所需的时间为$t$,返回的时间为$t_2$。现在，当光在从B跑向镜子的途中这段时间内，装置移动了一段距离$ut_1$,所以光必然用速度$c$走了$L+ut_1$的距离。我们也可以用$ct_1$来表示这段距离。这样，我们就有
\begin{equation*}
    ct_1=L+ut_1 ~\text{或}~ t_1=\frac{L}{(c-u)}
\end{equation*}
（这一结果也可以用这个观点明显看出，即光相对于装置的速度是$c-u$,所以这个时间是长度$L$除以$c-u$）。用类似方法可以算出时间$t_2$。在这段时间中B片前进了距离$ut_2$,所以光返回的距离是$L-u_2$。于是就有
\begin{equation*}
    ct_1=L+ut_1 ~\text{或}~ t_1=\frac{L}{(c-u)}
\end{equation*}
所以总的时间是
\begin{equation*}
    t_1+t_2=\frac{2Lc}{c^2-u^2}
\end{equation*}
为了便于以后对时间进行比较，我们把它写成
\begin{equation}
    \label{Eq:I:15:4}
    t_1+t_2=\frac{2L/c}{1-u^2/c^2}
\end{equation}

我们的第二部分，将是计算光从B到镜C的时间$t_3$。与前面一样，在$t_3$时间内，镜C向右移动了$ut_3$距离而到达位置C$'$;同时，光沿着一个直角三角形的斜边跑过距离$ct_3$,即BC$'$。对于这个三角形，我们有
\begin{equation*}
    (ct_3)^2=L^2+(ut3)^2
\end{equation*}
或
\begin{equation*}
    L^2=c^2t^2_3-u^2t^2_3=(c^2-u^2)t^2_3
\end{equation*}
由此可得
\begin{equation*}
    t_3=\frac{L}{\sqrt{c^2-u^2}}
\end{equation*}
从C$'$返回的路程与此相同，这可从图形的对称性上看出；所以返回用的时间也相同，因而整个时间是$2t_3$。稍微改变一下形式，我们可以写成
\begin{equation}
    \label{Eq:I:15:5}
    2t_3=\frac{2L}{\sqrt{c^2-u^2}}=\frac{2L}{c\sqrt{1-u^2/c^2}}
\end{equation}

现在我们就能比较两条光线所需的时间了。在（\ref{Eq:I:15:4}）和（\ref{Eq:I:15:5}）两式中，分子是相同的，它表示在装置静止不动的假定下光所取的时间。在分母中，除非$u$可以与$c$相比拟，不然$u^2/c^2$就很小。分母代表了由于装置运动而引起的时间上的修正。但必须注意，这些修正是不相同的——到达C来回所花的时间略小于到达E来回所花的时间，尽管两面镜子离B等距离，而我们所要做的就是要精确地把这个差别测量出来。

这里产生一个次要的技术性问题——假设两段长度$L$并不精确相等，怎么办？事实上，我们肯定不能使它们完全相等。在这种情况下，我们只要把装置转过90$^\circ$，使BC保持在运动的方向上，而BE则垂直于运动的方向。于是任何长度上的微小差别都变得不重要，我们要寻找的只是在装置转动时干涉条纹的移动。

在进行实验时，迈克耳逊和莫雷将仪器调整得使BE接近于平行地球的运动轨道（在白天和夜晚的一定时刻），地球的轨道速度约$18~mi \cdot s^{-1}$,那么，在一昼夜的某个时刻和一年之中的某些时候任何“以太漂移”都至少应有这么大。仪器的灵敏度足以观察到这种效应，但是，并没有发现时间上的差别——地球通过以太的速度无法被检测到。实验的结果说明不存在这种效应。

迈克耳逊-莫雷实验的结果令人迷惑不解和十分困扰。摆脱这个绝境的第一个有成效的观念是洛伦兹想到的。他提出，物体运动时会收缩，收缩只发生在运动方向上。这样，如果物体静止时长度为$L_0$，那么当它以速度$u$平行于其长度方向运动时，新的长度$L_\parallel$则为
\begin{equation}
    \label{Eq:I:15:6}
    L_\parallel = L_0 \sqrt{1-\frac{u^2}{c^2}}
\end{equation}
当将这个修正应用于迈克耳逊-莫雷干涉仪时，从B到C的长度没有改变，但是从B到E的长度缩短至$L\sqrt{1-u^2/c^2}$.因此，式（\ref{Eq:I:15:5}）没有改变，但式（\ref{Eq:I:15:4}）中的$L$必须按式（\ref{Eq:I:15:6}）而改变。这样做以后，我们得到
\begin{equation}
    \label{Eq:I:15:7}
    t_1 + t_2 = \frac{(2L/c)\sqrt{1-u^2/c^2}}{1-u^2/c^2}=\frac{2L/c}{\sqrt{1-u^2/c^2}}
\end{equation}
将此结果与式（\ref{Eq:I:15:5}）作比较，我们看到$t_1+t_2=2t_3$。所以如果仪器按刚才所说的方式缩短的话，我们就能够理解为什么迈克耳逊-莫雷实验根本没有测出这种效应。收缩假设虽然成功地解释了实验的负结果，但却也被认为这只是专门用来解释所遇到的困难而发明的，因而是过于牵强的。然而在用作发现以太风的许多其他实验中，也都出现了类似的困难。看来这是大自然反对人类的“阴谋”，它引进了某种新的因素来破坏每一个人原来以为应当测出速度$u$的实验现象。

人们终于认识到，正如庞加莱指出的那样，整个阴谋本身乃是一条自然法则！庞加莱于是便假定说，应当存在这么一条自然定律，即不可能用任何实验来发现以太风；也就是说，不可能测定绝对速度。

\section{时间的变换}

在检验收缩的概念是否与其他实验事实相协调时，结果发现假如时间也用方程组（\ref{Eq:I:15:3}）中第四个变换加以修正的话，那么每一件事就都很正确。这就是为什么从B到C再返回的过程中所花的这段时间$2t_3$。由那个在飞船上做实验的人算出时与由另一个观看飞船飞行的静止观察者算出的结果不一样的原因，对于在飞船上的人来说时间就是$2L/c$,但对另一个观察者来说，时间就是$(2L/c){\sqrt{1-u^2/c^2}}$（式\ref{Eq:I:15:5}）。换句话说，当一个外部观察者看到飞船中的人点燃雪茄烟时，所有的过程看来都比正常情况慢，而对飞船上这个人来说，每件事情都以正常方式进行。所以，不仅是长度需要缩短，时间测量仪器（“时钟”）显然也必须减慢。也就是说，从飞船上的人看来，钟记录下的时间过去一秒钟时，对于外面的人来说，它指示的是$1/\sqrt{1-u^2/c^2}s$。

必须看看钟的机构，并且注意它走动时的情况。由于这样做是相当困难的，我们将只选取一种非常简单的钟。我们所选择的是一个相当简陋的钟，但在原则上它是能工作的：这是一米尺，两端各有一面镜子，当我们在镜子间发出一个光信号时，光信号将一直来回传送着，当它往下跑时，每一次都会使这个钟“滴答”响一声，就像一个标准“滴答”钟一样。我们制作两个这样的钟，其长度完全相同，并把它们放在一起起动，使之同步。此后它们就会一直走得一样，因为它们的长度是相同的，光速也总是$c$。我们把其中一个钟让那个飞船上的人带着，他将尺的方向摆成垂直于运动的方向，这样，尺的长度不会发生变化。我们怎么知道垂直的长度不会发生变化呢？观察者和飞行者可以约好在擦过的一瞬间彼此在对方的$y$尺上刻下标记。根据对称性，两个标记必定具有同样的$y$和$y'$坐标，不然的话，当这两个人聚到一起比较结果时，有一个标记会高于或低于另一个标记，从而就能断定究竟谁在运动。

现在我们来看一下正在运动的钟上发生什么情况。在把钟带上飞船之前，那个人同意它是一个良好的、标准的钟，以后在飞船飞行中他也没有发现任何异常的现象，如果他发现了，就能知道自己在运动，因为假如有任何事情因运动而变化，他就能断定自己正在运动。但是相对性原理认为在匀速前进的系统中这是不可能的，可见不会产生任何变化。另一方面，当那个外部观察者在飞船的钟经过旁边时，他看到在镜子之间来回传送的光线“真正”取的是一条“之”字形的路径，因为尺总是横向运动的。我们在迈克耳逊-莫雷实验中已经分析过这种“之”字形运动。假定说，在一定时间内，尺朝前运动的距离正比于$u$，如图\ref{figure:光束经过的斜向路程示意图}所示，在同样时间内，光经过的距离正比于$c$，那么垂直距离就正比于$\sqrt{c^2-u^2}$。

这就是说，运动钟内光来回跑动的时间要长于静止钟内的时间。因此，对于运动钟来说，滴答声之间的表观时间以与图中所示的直角三角形的斜边同样的比例增长。（这就是我们方程式中平方根式的由来。）从图中也可以明显地看出$u$越大，运动的钟走得越慢。不仅这一类特定的钟会走得慢，只要相对论是正确的话，无论按什么原理工作的任何其他钟也都会慢下来，并且是以同样的比例慢下来——我们毋须进一步分析就可以说这句话，为什么如此呢？

为了回答这个问题，假定我们另外有两只做得完全相同的利用齿轮的钟，或者是根据放射性衰变或其他原理的钟。然后我们校准这些钟，使其与我们原先的钟严格同
\begin{wrapfigure}{l}{0.5\textwidth}
    \centering
    \includegraphics[width=0.48\textwidth]{Chapter15/光束经过的斜向路程示意图}
    \caption{\footnotesize （a）一个“光钟”静止在S$'$坐标系中；（b）一个同样的，相对于S坐标系运动的钟；（c）在正在运动的“光钟”中，光束经过的斜向路程示意图}
    \label{figure:光束经过的斜向路程示意图}
\end{wrapfigure}
步。当光在先前的两只钟上来回，并在到达时发出滴答声，新的钟也完成了某种循环，它们同时以双重符合的闪光、响声或其他信号表明这一点。在这两只钟中我们取一只放到飞船上去，和先前那只钟放在一起。也许这只钟不会变慢，而与那只静止的同样的钟走得一样，这样就与另一个运动钟不一致了。嘿！假如果真发生这种事，飞船上的人就能利用他的两只钟之间的不一致来确定飞船的速度，但是，我们已经假定这是不可能的。我们毋须知道任何有关会使新的钟产生这种效应的机理——我们只知道，不管理由如何，它都将同先前那只钟一样变慢。

现在，假如所有的运动钟都变慢，测量时间的任何方式都得出较慢的时间节拍，那么我们就得说：在一定的意义上飞船中的时间本身变慢了。在这里，所有的现象——人的脉搏，他的思维过程，他点燃雪茄烟的时间，以致他成长衰老的过程——所有这些事都必定以同样的比例变慢，因为他无法说出他正在运动。生物学家和医生有时会说：在飞船上癌的扩散所需的时间不一定会延长，但是从现代的物理学家来看这几乎是肯定的，不然，人们就能利用癌的扩散速度来确定飞船的速度了！

时间随着运动而变慢的一个非常有趣的例子与$\mu$子有关。这是一些经过平均寿命$2.2 \times 10^{-6}~s$后会自行蜕变的粒子。这种粒子可以在到达地球上的宇宙射线中找到，也可以在实验室里由人工制造。在射向地球时，其中有些粒子在半空中就蜕变了，其余的则在与物体碰撞而被留下之后才蜕变。很清楚，即使$\mu$子的速度同光一样快，在这样短的寿命内，它所走过的路程也不会超过600米以上。但是，虽然$\mu$子是在大气层的顶部，即大约10 km高的地方产生的，但是我们在大气层下面的实验室里的宇宙线中也确实找到了它们。这怎么可能呢？答案是：不同的$\mu$子各以不同的速度运动，其中有一些十分接近于光速。从它们本身的观点来看，它们只生存了大约2 μs,但是从我们的观点来看，它们的寿命要长得多——长到可以使它能到达地面。时间增长的因子已知是$1/\sqrt{1-u^2/c^2}$。对于各种速度的$\mu$子，人们非常精确地测量了它们的平均寿命，其数值与上述公式相当吻合。

我们并不知道为什么$\mu$子会蜕变，或者它的内部机理是什么，但是我们确实知道它的习性符合于相对论原理。这就是相对论原理的用途——它使我们甚至对那些知之不很多的东西作出预言。例如，在我们对于什么是使$\mu$子蜕变的原因获得一些概念之前，还是能够预言到，当它以光速的$9/10$的速度运动时，其所能生存的表观寿命为$(2.2 \times 10^{-6})\sqrt{1一9^2/10^2}~s$;我们的这种预言是成功的——这是一件好事情。

\section{洛伦兹收缩}

现在，我们回到洛伦兹变换式（\ref{Eq:I:15:3}）上来，并试图更好地理解坐标系$(x,~y,~z,~t)$与$(x',~y',~z',~t')$之间的关系。这些坐标系我们将分别称为S和S$'$系或乔和莫系。我们已经看到，第一个等式是建立在洛伦兹的沿$x$方向的收缩这个假设上的；我们如何来证明发生这样一种收缩呢？在迈克耳逊-莫雷实验中，我们根据相对论原理现在理解到，横臂BC不可能改变长度；然而，实验得到的结果为零就要求两个时间必须相等。所以，为了使实验得出零结果，看来纵臂BE必须缩短一个因子$\sqrt{1-u^2/c^2}$。从乔和莫所做的测量来说，这个收缩意味着什么呢？假定随同S$'$系沿$x$方向运动的莫是在用米尺测量某点的$x'$坐标。他用尺量了$x'$次，因此他认为这段距离是$x'$ m。但从S系的乔看来，莫却用了一根缩短了的尺，所以所测得的“真实”距离应当是$x'\sqrt{1-u^2/c^2}$ m。于是，当S$'$系离开S系跑过了距离$ut$时，S上的观察者将会说，在他的坐标系中测得的同一点的距离是
\begin{equation*}
    x=x'\sqrt{1-\frac{u^2}{c^2}}+ut~\text{或}~x'=\frac{x-ut}{\sqrt{1-u^2/c^2}}
\end{equation*}
这就是洛伦兹变换的第一个等式。

\section{同时性}

同样的情况表明，由于时间尺度上的不同，分母的表示式也被引进到洛伦兹变换的第4个等式中。这个等式中最有趣的一项是分子中的$ux/c^2$项，因为它是全新的，而且是未曾预料到的。那么这究竟意味着什么呢？如果我们仔细地来看一下这个情况，我们可以发现，发生在不同地点的两个事件，在S$'$中的莫看来发生于同时，但在S中的乔看来，它们并不发生于同时。如果一个事件在$x_1$处发生于时间$t_0$，另一个则在$x_2$发生于时间$t_0$（同一时刻），那么我们发现，两个相应的时间$t_1$与$t_2$相差一个量
\begin{equation*}
    t'_2-t'_1 = \frac{u(x_1-x_2)/c^2}{\sqrt{1-u^2/c^2}}
\end{equation*}
这种情况称为“异地同时性的破坏”。为了使这个概念稍为清楚一些，让我们考虑下面一个实验。

假设有一个人在运动的宇宙飞船上（系统S$'$）的两端各放置一只钟，并且想弄明白这两只钟是否已对准。怎样使这两只钟对准呢？有许多方法：一个方法只需要很少一点计算，这就是首先精确地确定两只钟之间的中点。然后从这个位置上发出一个光信号，这个光信号将以同样的速度沿两条路径传播，而且非常清楚将同时到达两只钟。信号的这种同时到达性可以用来把钟对准。我们假定S$'$中的人是用这种特殊方法对准他的钟的。我们再看一下S系统中的一个观察者是否会同意这两只钟已经对准。S$'$系统中的人相信这一点，因为他不知道他正在运动。但是S系统中的人则推论说，由于宇宙飞船向前运动，飞船前端的一只钟将离开光信号而去，因此为了追到它，光必须走过大于一半距离的路程；但后面的一只钟却迎着光信号而去，所以这段距离就较短。因此，信号会先到达后面一只钟，虽然S$'$中的人认为信号是同时到达的。因此我们看到，当宇宙飞船中的人认为两个地方的时间是同时的时候，在他的坐标系中的两个相等的$t'$值，必须对应于另一个坐标系中的两个不同的$t$值！

\section{思维矢量}

让我们再看看，从洛伦兹变换中还可以发现一些什么。有趣的是，可以注意到$x$项与$t$项之间的变换在形式上与我们在第11章中对于坐标系的转动曾研究过的$x$项和$y$项的变换非常相似。在那里我们有
\begin{equation}
    \label{Eq:I:15:8}
    \begin{aligned}
        x'&=x\cos\theta + y\sin\theta \\
        y'&=y\cos\theta - x\sin\theta \\
    \end{aligned}
\end{equation}
可见新的$x'$项混合了原来的$x$与$y$，新的$y'$项也混合了原来的$y$与$x$；与此相似，在洛伦兹变换中，我们发现新的$x'$是$x$与$t$的混合项，新的$t'$是$t$与$x$的混合项。这样，洛伦兹变换就类同于一种转动，不过这是一种在空时中的“转动”。这看来是一个奇怪的概念。这种与转动的类比，可以由下列量的计算而得到核实
\begin{equation}
    \label{Eq:I:15:9}
    x'^2 + y'^2 + z'^2 - c^2t'^2 = x^2 + y^2 + z^2 - c^2t^2
\end{equation}
在这个等式中，每一边的前三项在三维几何中所代表的是一点与原点之间的距离（一个球面）的平方。这个平方在坐标轴的转动下保持不变（不变量）。与此相似，式（\ref{Eq:I:15:9}）表明，存在着包含时间在内的某一种组合，它在洛伦兹变换下也是不变的。这样，与转动的类比就完全了，而且这是这样一种类比，即矢量，也就是其中包含与坐标和时间以同样方式转换的“分量”的那些量，对于相对论也是有用的。

于是我们试图把矢量的观念加以扩展，使其包括时间的分量，而在此以前，我们认为它只有空间分量。这就是说，我们期望将有一种具有四个分量的矢量，其中三个同一般矢量的分量一样，而与这些分量一起还加上第四个分量，它是时间部分的类比项。

这个概念将在后面几章中继续加以分析，在那里我们将看到，如果把前一节中的观点应用于动量时，那么变换将给予我们三个空间部分，它们如同通常的动量分量一样，另外还有第四个分量，也就是时间部分，那正是能量。

\section{相对论力学}

我们现在已经为更一般地研究在洛伦兹变换下力学定律将采取什么形式作好了准备。[到目前为止，我们说明了长度和时间如何变化，但没有说明我们是如何得到$m$的修正公式（式\ref{Eq:I:15:1}）的。我们将在下一章中来加以说明。]为了看出爱因斯坦对牛顿力学的质量$m$进行修正的重要意义，我们从牛顿第二定律出发，即力是动量的变化率为
\begin{equation*}
    \boldsymbol{F} = \frac{d(m\boldsymbol{v})}{dt}
\end{equation*}
动量仍然是$m\boldsymbol{v}$，但当我们用新的$m$时，它就变为
\begin{equation}
    \label{Eq:I:15:10}
    \boldsymbol{P} = m\boldsymbol{v} = \frac{m_0\boldsymbol{v}}{\sqrt{1-v^2/c^2}}
\end{equation}
这就是爱因斯坦对牛顿定律的修正。在这种修正下，如果作用和反作用仍然相等（这不一定指每个时刻，而就最终结果来说是相等），那么动量守恒仍像以前一样成立，但是守恒的量，不再是原来的具有不变质量的 $m\boldsymbol{v}$，而是如式（\ref{Eq:I:15:10}）所表示的，具有经过修正的质量的量，如果在动量公式中考虑到这种变化，那么动量守恒定律仍然有效。

现在我们来看看，动量如何随速度而变化。在牛顿力学中，它正比于速度，而且按照式（\ref{Eq:I:15:10}）在速度与光速相比甚小的一个相当大的范围内，它在相对论力学中近乎与之相同，也正比于速度，因为平方根这个因子与1相差甚微。但是当$v$几乎等于$c$时，分母的平方根趋向于零，因此动量趋向于无穷大。

如果有一个恒力作用在一个物体上很长时间，那么会出现什么情况？牛顿力学认为，物体将不断获得速度，直到它的运动超过光速。但是在相对论力学中，这是不可能的。在相对论中，物体不断得到的不是速度，而是动量。动量可以因为质量在不断增加而持续增大。经过一定时间后，实际上已不存在那种在速度变化含义上的加速运动，但是动量却继续在增加。自然，如果一个力只使物体的速度产生非常小的变化，我们就说这个物体具有很大的惯性。这正是我们的相对论质量公式所指出的[见式（\ref{Eq:I:15:10}）]：当v大到接近于c时，惯性是非常大的。作为这种效应的一个例子，我们举出，在加利福尼亚理工学院所使用的同步加速器中，为了要偏转高速电子，所需的磁场的强度要比依据牛顿定律所预言的大2000倍。换句话说，同步加速器中的电子的质量为它们正常质量的2000倍，就如同一个质子的质量那么大！ $m$是$m_0$的2000倍，意味着（$1—v^2/c^2$）必定为$1/4000000$，也就是说$v^2/c^2$与1的差别只是四百万分之一。或者说$v$与$c$的差别只有$c$的八百万分之一。所以电子的速度非常接近于光的速度。如果电子与光同时开始从这个加速器射到邻近一个实验室（约700 ft远），那么谁先到达呢？当然是光，因为光总是跑得更快一些\footnote{在同可见光的比赛中，由于空气的折射率，实际上电子将赢得胜利。但$\gamma$射线会稳操胜券}。但是早到多少时间呢？回答起来太麻烦了，我们还是说光所超前的路程有多少：约为$1/1000$ in，或者说一张纸的厚度的$1/4$！当电子跑得这样快时，它们的质量是非常巨大的，然而它们的速度不会超过光的速度。

现在我们再看看质量的相对论效应所具有的其他一些结果。考虑在一个小的容器中气体分子的运动。当气体被加热时，分子的速度就增加，因此，它的质量也会增加，而气体变重了。当速度较小时，表示质量增加的一个近似公式，可以利用二项式定理把
\begin{equation*}
    \frac{m_0}{\sqrt{1-v^2/c^2}}=m_0(1-\frac{v^2}{c^2})^{-\frac{1}{2}}
\end{equation*}
展开为幂级数而得到。我们得到
\begin{equation*}
    m_0(1-\frac{v^2}{c^2})^{-\frac{1}{2}}=m_0(1+\frac{1}{2} \frac{v^2}{c^2}+\frac{3}{8} \frac{v^4}{c^4}+\cdots)
\end{equation*}
从这个表示式可以清楚看出，当$v$较小时，级数收敛得很快，在前二项或前三项之后的各项可以忽略不计。因而我们可以写为
\begin{equation}
    \label{Eq:I:15:11}
    m \approx m_0+\frac{1}{2}m_0v^2(\frac{1}{c^2})
\end{equation}
其中，右端的第二项表示由分子的速度而来的质量的增加。由于温度升高时，$v^2$与之成正比地增加，所以我们可以说，质量的增加正比于温度的增加。由于$m_0v^2/2$在原来的牛顿含义中是动能，所以我们也可以说，整个气体的质量的增加，等于动能增加的量除以$c^2$，或者说$\Delta m = \Delta (\text{K.E.})/c^2$。

\section{质能相当性}

上面的观察给了爱因斯坦一个启发，使他想到，如果我们说物体的质量等于该物体总的能量含量除以$c^2$，那么一个物体的质量就可以表示得比式（\ref{Eq:I:15:1}）更为简单。如果式（\ref{Eq:I:15:11}）乘以$c^2$，则结果为
\begin{equation}
    \label{Eq:I:15:12}
    mc^2=m_0c^2+\frac{1}{2}m_0v^2+\cdots
\end{equation}
这里，左端的一项表示一个物体的总能量，右端的后面一项可以认为是通常的动能。爱因斯坦把很大的常数项$m_0c^2$解释为该物体总能量的一部分，是一种通常称为“静能”的内在能量。

让我们跟着爱因斯坦来探究物体的能量总是等于$mc^2$这个假设会得出一些什么结论。作为一个有趣的结果，我们将找出表示质量随速度而变化的公式（\ref{Eq:I:15:1}），而迄今为止我们只是把它作为一个假设来看待。我们从处于静止状态的一个物体出发，其能量为$m_0c^2$。然后对这个物体施加一个力。这个力使物体开始运动，并给予它动能；因此，由于能量增加，质量也增加——这已包含在原来的假设之中。只要力继续作用在物体上，能量和质量两者都会继续增加。我们已经看到（第13章），能量对时间的变化率等于力乘以速度，或
\begin{equation}
    \label{Eq:I:15:13}
    \frac{dE}{dt}=\boldsymbol{F} \cdot \boldsymbol{v}
\end{equation}
我们还看到[第9章，式（\ref{Eq:I:9:1}）] $F=d(mv)/dt$。当这些关系与$E$的定义结合在一起，式（\ref{Eq:I:15:13}）就变为
\begin{equation}
    \label{Eq:I:15:14}
    \frac{d(mc^2)}{dt}=\boldsymbol{v} \cdot \frac{d(mv)}{dt}
\end{equation}
我们希望解这个关于$m$的方程。为此，我们先用一点数学技巧，即在式子两端各乘以$2m$，这就把方程变为
\begin{equation}
    \label{Eq:I:15:15}
    c^2(2m)\frac{dm}{dt}=2mv \frac{d(mv)}{dt}
\end{equation}
我们要除去微商，这可以在两端用积分来做到。可以看出量$(2m)\cdot dm/dt$是$m^2$的时间微商，而$(2m\boldsymbol{v}) \cdot d(m\boldsymbol{v})/dt$则是$(mv)^2$的时间微商。这样，等式（\ref{Eq:I:15:15}）就等于
\begin{equation}
    \label{Eq:I:15:16}
    \frac{c^2d(m^2)}{dt}=\frac{d(m^2v^2)}{dt}
\end{equation}
假如两个量的微商相等，它们本身最多只差一个常数，比如说$C$。这就使我们能写成
\begin{equation}
    \label{Eq:I:15:17}
    m^2c^2=m^2v^2+C
\end{equation}
现在必须把这个常数$C$定义得更清楚一点。由于式（\ref{Eq:I:15:17}）必须对所有的速度都成立，所以我们可以选择$v=0$这个特殊情况，并且说这时的质量是$m_0$。将这些值代入等式（\ref{Eq:I:15:17}），得出
\begin{equation*}
    m_0^2c^2=0+C
\end{equation*}
现在我们可以把这个$C$值代入式（\ref{Eq:I:15:17}）中，于是得到
\begin{equation}
    \label{Eq:I:15:18}
    m^2c^2=m^2v^2+m_0^2c^2
\end{equation}
除以$c^2$，再把各项整理一下后得
\begin{equation*}
    m^2(1-\frac{v^2}{c^2})=m_0^2
\end{equation*}
由此便得
\begin{equation*}
    m=\frac{m_0}{\sqrt{1-v^2/c^2}}
\end{equation*}
这就是式（\ref{Eq:I:15:1}），也正是为了使式（\ref{Eq:I:15:12}）中质量与能量之间相符合所必要的一个公式。

通常说来，能量的这种变化只表示质量上极其微小的变化，因为平时我们不可能从一定量的物质中产生很多能量；但是，在原子弹爆炸中，如果其能量相当于20000吨TNT炸药，就可以知道爆炸后的尘埃将比反应材料的原有质量轻1g，因为，按照关系式$\Delta E=\Delta (mc^2)$，所释放的能量相当于1g的质量。这种质能等价性的理论已为由物质没而完全转化为能量的实验出色地证实了：一个负电子与一个正电子在静止时质量各为$m_0$，当它们碰到一起时，会蜕变成两束$\gamma$射线，测得各带有$m_0c^2$的能量。这个实验为确定与粒子的静止质量的存在相关的能量提供了一个直接的方法。